The operation inventory in Tables~\ref{tab:audit-counts} and
\ref{tab:audit-counts-additional} distinguishes retained raw counters from
exact code-based derivations. One likelihood evaluation ($L$) concerns one
particle state, including a four-dimensional state in QGPF-like; vectorized
calls count their particle entries. $C$ counts generated candidate states or
binary subsets submitted to scoring. $P$ preserves each implementation's
historical counter: a candidate attempt for QM-HPF-like and the budget audits,
but an offspring search episode for QGPF-like, which contains four candidate
generations. $R_b$, $X$, and $M_b$ denote classical bit rotation routine calls,
arithmetic crossover proposals, and Bernoulli bit observations, respectively.
QM bit routine counts include zero-angle/no-change branches. $A_b$ counts
persistent probability-entry updates, excluding initialization; the local
probability computations in QM-HPF-like do not maintain such a persistent
state. $F$ and $D$ retain the EC-QPF search-fitness and candidate-distance
counter conventions, excluding final diagnostic recomputations. The exact
classical comparator's $F$ is an argmin scan over singleton entries, rather
than one exponentiation per entry. $G$ counts a low-particle generation for
QM-HPF-like, an offspring generation for QGPF-like, and a full swarm generation
for the EC-QPF module.

Branch-level totals absent from the original logs remain
\textit{unavailable}; supplementary seed-0 instrumentation does not replace
missing 30-seed means. A transparent-wrapper replay recovered the historical
native and QGPF-like seed-0 estimates and recorded counters without consuming
additional random numbers. All bit operations are classical simulations, and
these counts concern the audited implementations and selection module. The
columns overlap and their units differ: an EC-QPF generated mask is also a
fitness evaluation, while a QGPF rotation also updates a probability entry.
Consequently, these counts are neither additive operation equivalents nor
evidence of physical quantum acceleration. Equal likelihood budgets do not
imply equal wall-clock costs.
